% !TeX program = xelatex
\documentclass{UAmathtalk}
\usepackage{setspace}
\renewcommand{\when}{12:00 p.m.}
\renewcommand{\where}{Hudson 0110}
%\renewcommand{\where}{Hudson 0110.\\ \href{https://albany.zoom.us/j/81219478289?pwd=N5fy7vxKBXqlQ4mAq4KaeaAM5iAHi4.1}{Zoom ID: 812\,1947\,8289 - Passcode: 233444}}
\author{ Emilio Minichiello }
\address{CUNY CityTech}
\urladdr{https://www.emiliominichiello.com/}
\title{GRAPH HOMOTOPY THEORY}
\date{Friday, October 16th, 2026}


\begin{document}

\maketitle
%\renewcommand*{\abstractsize}{\normalsize}
\renewcommand*{\abstractsize}{\fontsize{13.5}{16.2}\selectfont}
\begin{abstract}
    {
The area of topological combinatorics has grown immensely in recent years. There are now many ways to construct a space from a discrete structure like a graph, matroid or simplicial complex and extract some combinatorial information about the discrete structure via the topology of the resulting space. In this talk I will go over my paper “Thomason-Type Model Structures on Simplicial Complexes and Graphs,” which goes over a particular way to do this, called reflexive X-homotopy theory. I show that by probing a graph G with reflexive complete graphs we obtain a space Sing G, whose construction factors through the clique complex of G. Matsushita showed that one can transfer a model structure from simplicial sets to loop graphs via this construction, and I proved that his construction actually factors through reflexive graphs and simplicial complexes, and one can transfer model structures to these categories as well. Furthermore by factoring Matsushita’s construction in this way, we obtain a significant amount of information on how to actually compute Sing G, as well as obtaining nice properties of the model categories involved, like properness and (co)fibrancy conditions for certain classes of graphs. The talk will be kept at an elementary level, with plenty of pictures and examples.
    }
\end{abstract}
\end{document}


